\section{相関分析 — 関係を測るという思想と実践}

\subsection{問い（Question）}
日常の中で, 私たちは「何かと何かが一緒に変化している」感覚を頻繁に持つ.  
例えば, 気温が高い日にアイスクリームの売上が増える, 広告費を増やすと売上も伸びる, 睡眠時間が長い人ほど集中力が高いといった観察である.  
これらは経験則としては有用だが, 科学的に正しいかどうかは別問題である.  

ここでの問いは二つある.  
第一に, 観察された「一緒に変化しているように見える」現象は, 本当に統計的な関係を持っているのか.  
第二に, その関係はどの程度の強さで, どのような方向性を持つのか.  

この問いに答えるためには, 人間の主観や印象に左右されない, 客観的で再現可能な方法が必要になる.  
相関分析は, 二つの変数の間の関係性を一つの数値で表現する方法であり, 本章ではその思想・理論・実践・限界を体系的に学ぶ.

\subsection{素朴な方法と限界（Naive Approach \& Limitations）}
最初に思いつくのは, 二つの変数を並べて比較する方法である.  
例えば, 気温とアイス売上の折れ線グラフを描き, 山や谷のタイミングが似ていれば「関係がありそうだ」と感じる.  
散布図にして右上がりなら正の関係, 右下がりなら負の関係という読み取りもできる.  

しかし, この方法には主観的判断という重大な弱点がある.  
目の錯覚や事前の期待によって, 関係が弱くても強く見えたり, 逆に強い関係があっても見逃したりする.  
さらに, 軸のスケールや表示期間を変えると, 見た目の印象が大きく変わってしまう.  

もう一つの限界は, 見た目では関係の強さを定量的に比較できないことだ.  
例えば, 二つの散布図がどちらも右上がりに見える場合, どちらの関係がより強いのかは, 数値化しないと判断できない.  

このような理由から, 直感的観察だけに頼らず, 数式に基づく客観的な指標が不可欠となる.

\subsection{理論の最小核（Theory Kernel）}
相関分析の出発点は, 「変数同士が同じ方向に動くかどうか」という直感的な問いである.  
これを数学的に定義するため, 平均からのずれ（偏差）を利用する.  

変数 $X$ の平均を $E[X]$, 偏差を $(X - E[X])$ とする.  
二つの変数 $X, Y$ の偏差を掛け合わせた $(X - E[X])(Y - E[Y])$ は, 同方向に動くとき正, 逆方向のとき負になる.  
これを平均したものが共分散である:
\[
\mathrm{cov}(X,Y) = E[(X - E[X])(Y - E[Y])].
\]

共分散が正なら「増えると増える」の関係, 負なら「増えると減る」の関係, 0なら線形関係がほぼないことを示す.  

しかし, 共分散は変数の単位やスケールに依存する.  
例えば, 身長をcm単位で測るかm単位で測るかで, 共分散の値は100倍変わってしまう.  
これでは異なる変数ペア間の比較ができない.  

そこで, 各変数の標準偏差で割ってスケールを除去した指標が, ピアソンの相関係数である:
\[
\rho_{XY} = \frac{\mathrm{cov}(X,Y)}{\sigma_X \sigma_Y}, \quad \sigma_X = \sqrt{V(X)}, \quad \sigma_Y = \sqrt{V(Y)}.
\]

さらに, 標準化変数
\[
Z_X = \frac{X - E[X]}{\sigma_X}, \quad Z_Y = \frac{Y - E[Y]}{\sigma_Y}
\]
を用いれば,
\[
\rho_{XY} = E[Z_X Z_Y]
\]
となる.  
これは「単位を揃えた二変数の共分散」であり, 相関係数の概念的本質を端的に示す.

\subsection{手法（Method）}
標本データに基づく相関係数の計算手順は以下の通りである.

1. 標本平均 $\bar{X}, \bar{Y}$ を計算する.  
2. 各観測値の偏差 $(X_i - \bar{X}), (Y_i - \bar{Y})$ を求める.  
3. 偏差同士を掛け合わせ, それらを合計する.  
4. 合計値を $(n-1)$ で割って標本共分散を求める.  
5. 標本標準偏差 $s_X, s_Y$ を求める.  
6. 標本相関係数 $r_{XY}$ を
\[
r_{XY} = \frac{\mathrm{cov}(X,Y)}{s_X s_Y}
\]
で計算する.

計算過程で, 偏差の符号と大きさがどのように積に反映されるかを確認すると, 相関係数の意味が直感的に理解できる.

\subsection{実践（Practice）}
仮想データとして, 10カ月間の広告費（万円）と売上（万円）を考える.  
散布図を描くと, 明らかに右上がりの帯状に点が並び, 計算すると $r = 0.87$ で強い正の相関が得られる.  

しかし, この情報だけで「広告費が売上を増やす」と結論するのは危険である.  
例えば, 季節要因やイベントなど, 第三の要因が両者に影響している可能性がある.  
これは因果の混同であり, 擬似相関の典型例である.  

さらに, もし1つの月だけ極端な広告費と売上を記録していれば, その1点が相関係数を大きく変えてしまうこともある.  
外れ値の影響を防ぐため, 相関分析の前には必ず散布図を確認し, データの分布を目視で把握するべきである.

\subsection{結果と次の一手（Output \& Next）}
本章では, 相関係数の理論的背景と計算方法, 実践例とその注意点を学んだ.  
次の段階として, 非線形の関係や順位相関係数, さらには因果推論との違いを学ぶことで, より多様なデータ構造に対応できるようになる.


\subsection*{補足：相関分析の罠}
相関係数は二つの変数の関係を数量化する強力な道具であるが, その解釈を誤ると致命的な結論を導いてしまう.  
常に警戒すべき三つの典型的な落とし穴を, 簡単な事例と共に見ておこう.

\subsubsection*{1. 因果の混同と疑似相関 — アイスクリームと水難事故}
夏になると, 「アイスクリームの売上」と「水難事故の件数」はともに強い正の相関を示す.  
だからといって「アイスを食べると溺れる」と考えるのは明らかな誤りだ.  
両者を同時に動かしているのは「気温の上昇」という第三の要因である.  
このように, 二つの変数の背後に共通の原因（交絡因子）が潜んでいる場合, 見かけ上の相関は本質的な関係を意味しない.  
この錯覚を避けるには, 可能性のある第三の要因を列挙し, データ収集や分析でその影響を取り除くことが重要である.

\subsubsection*{2. 外れ値の影響 — たった一つの点が全体を歪める}
普段は無相関に見えるデータでも, 一つだけ極端に大きな値のペア（外れ値）が存在すると,  
その一点に引きずられて, 見かけ上は高い相関係数が計算されることがある.  
逆に, 強い相関関係を打ち消してしまう外れ値も存在する.  
例えば, 座標平面にほぼ一直線に並んだ点群があっても, 離れた場所に1点だけ異常な点を置けば, 相関係数は大きく変動する.  
このため, 相関係数を計算する前には必ず散布図を描き, 「おかしな点」がないかを目視で確認する習慣を持つことが欠かせない.

\subsubsection*{3. シンプソンのパラドックス — 部分と全体の矛盾}
統計学における最も直感に反する現象の一つに, シンプソンのパラドックスがある.  
例えば, ある大学の学部Aと学部Bの合格率を調べると, 男女別に見ればどちらの学部でも女性の合格率の方が高い（女性に有利な相関）が,  
大学全体で合計すると, 男性の合格率の方が高くなる（男性に有利な相関）という事態が起こりうる.  
この逆転は, 各学部の志願者数に男女で大きな差があることが原因である.  
つまり, 集計の仕方によって相関の符号すら変わってしまうのだ.  
このパラドックスは, データを集計する際には背後にあるグループ構造や分布の偏りを無視してはならないという重要な教訓を与えてくれる.

  
  



  % \section{相関分析}

  % 「ついてない日は不幸が重なる気がする」  
  % 「顔がよい人は, なんだかんだでモテる」  
  % 「気圧が低い日は, 頭が痛くなりやすい」  
  % こんなふうに，私たちは日常の中で「何かと何かが関係している」と感じることがよくある．

  % こうした「関係」はしばしば，私たちの予測や判断の根拠となる．  
  % しかし，その関係は本当に存在するのだろうか？ どのくらい強い関係なのか？ どちらが原因で，どちらが結果なのか？  
  % このような問いに，データをもとに向き合おうとするのが，\textbf{相関分析}である．

  % 本章では，変数同士の関係を「相関」という観点から捉えるための基本的な考え方と道具立てを紹介する．  
  % 数式的な定義だけでなく，散布図などを通じて「関係を見る」という行為そのものにも注目し，  
  % 統計的な視点がどのように直感的理解と結びついていくかを順にたどっていく．


  % \subsection{データどうしの関係を見るための基本的な考え方}
  % まず，「データどうしの関係を見る」とはどういうことかを考えてみよう．  
  % たとえば，天気と売上の関係を調べたいと思ったとき，それぞれのデータを折れ線グラフにして重ねてみると，  
  % 両者が似たような動きをしているように見えることがある．

  % 「晴れの日に売上が伸びて，雨の日に落ちる」といった傾向が見えたとき，私たちは「関係がある」と感じるだろう．  
  % このように，2つの量がある意味で「一緒に変化している」ように見えるとき，そこにはデータの\textbf{相関}があるかもしれない．

  % では，この「一緒に動く感じ」を，ただの印象だけでなく，数値として表すにはどうすればよいのだろうか？  
  % 実際に何をもって「一緒に動く」と言えるのかは，すぐにははっきりしない．

  % たとえば，$X$ の値が大きくなるときに $Y$ の値も大きくなる，あるいは逆に小さくなる，というように，  
  % 変数同士の変化の向きが何らかの仕方で連動していると感じるとき，私たちは「一緒に動いている」と言いたくなる．

  % このような関係を定量的に捉えるには，「変化の方向」というものを何らかのかたちで表現する必要がある．  
  % そこで手がかりになるのが，それぞれの変数が持つ\textbf{平均値}である．  
  % 平均はその変数の「中心」と考えられるため，その値を基準として，各データ点がそこからどのように離れているか（＝ずれているか）を見ることで，  
  % 「今この値は大きいのか小さいのか」という情報を取り出すことができる．

  % 平均からのずれに注目すれば，「$X$ が平均より大きいとき，$Y$ も大きいか？」といった問いを形式化できる．  
  % このずれを使って，ふたつの変数の変化の向きを比較することができるのである．

  % この考えに基づいて定義されるのが，共分散である：
  % \[
  % \mathrm{cov}(X,Y) = E[(X - E[X])(Y - E[Y])].
  % \]

  % この式では，$X$ と $Y$ のそれぞれの値が平均からどのくらい離れているか（ずれ）を掛け合わせ，  
  % その平均をとっている．

  % このときの掛け算は，「変化の向きの一致」に関する情報を取り出すために機能している．  
  % 符号の組み合わせを考えてみれば，その理由がわかる：

  % \begin{itemize}
  %   \item 両方のずれが正（平均より大きい） → 積は正
  %   \item 両方のずれが負（平均より小さい） → 積は正
  %   \item 一方が正で他方が負 → 積は負
  % \end{itemize}

  % このように，$(X - E[X])(Y - E[Y])$ は，それぞれの変数が平均からどのように“ずれているか”という方向の情報を掛け合わせることで，  
  % 2つの変数の“ずれの向きの一致”を測る役割を果たしている
  % （ベクトルの内積のように考えれば，「向きが揃っていれば正，逆なら負」という直感にも通じる）.
  % いろいろなデータ点についてこの積を平均すれば，全体として「一緒に動く傾向」があるかどうかが測れることになる．

  % この定義は，分散の式
  % \[
  % V(X) = E[(X - E[X])^2]
  % \]
  % と並べて見ると構造的によく似ている．  
  % 分散が，1つの変数がどのくらいばらついているかを測っているのに対し，  
  % 共分散は，2つの変数がどのように同時に変化しているかを測っている．

  % 共分散が正の値であれば，$X$ が平均より大きいときに $Y$ も大きい傾向があることを意味し，  
  % 負の値であれば，一方が大きいときに他方は小さくなる傾向があることを示す．  
  % 関係がなければ，正と負の積が打ち消し合い，共分散はゼロに近づく．

  % たとえば，スマートフォンの利用時間と利用している人の自尊感情には，負の関係があるという調査結果がある．  
  % スマホの利用時間が長い人ほど，自分に対する肯定的な気持ちが低い傾向が見られるとすれば，  
  % この2つの変数は逆方向に変化していることになり，共分散は負の値をとる．

  % 一方で，ある音楽アプリの「再生回数」と「お気に入り登録数」のように，  
  % どちらもユーザーの関心を表す指標であり，一緒に増える傾向がある場合には，  
  % 両者の共分散は正になると考えられるのだ．


  % \subsection{ピアソンの相関係数}
  % 共分散は，ふたつの変数が「同じ方向に動いているかどうか」を捉えるための指標であった．  
  % しかし，それだけでは「どのくらい強く関係しているか」を，他の変数の組と比べることが難しいという問題がある．

  % たとえば，ある企業が，「広告費と売上の関係」と「社員数と売上の関係」のどちらが強いかを知りたいとしよう．  
  % それぞれについて共分散を計算してみたところ，前者は 1200，後者は 85 だったとする．  
  % しかし，広告費は数百万円単位で動き，社員数は百人単位でしか変化しないとすれば，  
  % この2つの共分散をそのまま比較して，「広告の方が売上に関係していそうだ」と結論づけることはできない．

  % 変数のスケールが異なる場合，\textbf{共分散の値は大きさの比較に意味を持たない}．  
  % 同じ売上への影響を見ているにもかかわらず，単位やばらつきによって共分散の値そのものが変わってしまう．

  % この問題を解決するために，共分散をそれぞれの変数のばらつき（分散）で割って，  
  % \textbf{スケールの影響を取り除いた}ものが定義される．  
  % これが，\textbf{ピアソンの相関係数}である：

  % \[
  % \rho_{XY} = \frac{\mathrm{cov}(X,Y)}{\sqrt{V(X)} \cdot \sqrt{V(Y)}}.
  % \]

  % 相関係数 $\rho_{XY}$ は，$-1$ から $1$ の間の値をとる．  
  % 値が $1$ に近いときは，「完全な右上がり」の関係（直線的で，$X$ が大きくなると $Y$ も必ず大きくなる）を示し，  
  % $-1$ に近いときは，「完全な右下がり」の関係を示す．  
  % 0 に近ければ，線形な関係が見られないことを意味する．

  % このように，相関係数は，共分散と同じように「関係の方向」を捉えながら，  
  % 異なる単位やスケールの変数どうしでも，関係の強さを比較できるようになっている．

  % \begin{remark}
  %   ピアソン相関係数は, 線形変換（例：$aX+b$, $cY+d$; $a,c>0$）をしても値が変わらない．

  %   たとえば，ある人のスマートフォンのスクリーンタイム $X$ と睡眠時間 $Y$ の間に相関があるとする．  
  %   スクリーンタイムを「時間」ではなく「分」で測り直したり，睡眠時間を「実測値」ではなく「平均からの差」として記録し直したとしても，  
  %   変数の関係の強さや向き（右上がりか右下がりか）は変わらない．  
  %   このように，相関係数は変数の単位や尺度の変更に左右されず，\textbf{関係そのものの構造}を捉えることができる．
  % \end{remark}

  % \begin{remark}
  %   相関係数は，次のようにも理解することができる．  
  %   変数 $X$ と $Y$ を，それぞれ平均 $0$，分散 $1$ に変換したうえで，共分散をとるという見方である．  
  %   このような変換を，\textbf{標準化}という：

  %   \[
  %   Z_X = \frac{X - E[X]}{\sqrt{V(X)}}, \quad Z_Y = \frac{Y - E[Y]}{\sqrt{V(Y)}}.
  %   \]

  %   このとき，ピアソンの相関係数は，
  %   \[
  %   \rho_{XY} = E[Z_X Z_Y]
  %   \]
  %   と書くことができる．

  %   この見方によれば，相関係数は「標準化された $X$ と $Y$ の共分散」にほかならない．  
  %   つまり，「単位や大きさの異なる変数を，基準を揃えたうえで関係を見る」という発想が，数式の形にも現れている．
  % \end{remark}

  % \begin{remark}
  %   複数の変数どうしの関係を同時に見たいときには, 相関係数を表のようにまとめた「相関行列」というものが使われることがある．

  %   たとえば, $X_1$（運動時間）, $X_2$（睡眠時間）, $X_3$（スマホ利用時間）という3つの変数があったとすると, それぞれの組み合わせについて相関係数を計算し，次のような形に整理することができる：

  %   \[
  %   \begin{pmatrix}
  %   1 & \rho_{12} & \rho_{13} \\
  %   \rho_{21} & 1 & \rho_{23} \\
  %   \rho_{31} & \rho_{32} & 1 \\
  %   \end{pmatrix}
  %   \]

  %   これは，見た目には少し複雑だが，「関係が強い組み合わせはどれか」を一目で見渡すことができるという利点がある．  
  %   本書では詳細には立ち入らないが，「相関係数が複数あっても考え方は同じである」ことを確認するために，このような表現もあるということを紹介しておく．
  % \end{remark}

  % \subsection{相関係数と散布図の関係}

  % 相関係数は，2つの変数の間にどれくらい明確な「方向性のある関係」があるかを数値で表すものだった．  
  % この関係は，グラフとして描いたときの点の並び方（散布図）として視覚的に感じ取ることもできる．

  % たとえば，$X$ を横軸に，$Y$ を縦軸にとって，いくつかのデータ点を散布図にすると，  
  % 次のような違いが見られる：

  % \begin{itemize}
  %   \item 相関係数 $\rho = 1$ のとき，すべての点がちょうど1本の直線（右上がり）上に並ぶ．
  %   \item $\rho = 0.8$ のとき，多くの点が右上がりの帯状に分布し，「上がれば上がる」傾向が見られる．
  %   \item $\rho = 0$ のとき，典型的には点の並びに明確な方向性は見られず，雲のように広がる．  
  %         （ただし，見た目には何らかの関係がありそうに見えても，相関係数が 0 に近くなることもあるので注意が必要である．）
  %   \item $\rho = -0.8$ のとき，右下がりの帯状に広がり，「一方が上がれば他方が下がる」傾向が見られる．
  %   \item $\rho = -1$ のとき，すべての点が右下がりの直線上に並ぶ．
  % \end{itemize}

  % \vspace{0.5em}

  % \noindent
  % \textbf{[図挿入：$\rho = 1$ の散布図]}  
  % \textbf{[図挿入：$\rho = 0.8$ の散布図]}  
  % \textbf{[図挿入：$\rho = 0$ の散布図]}  
  % \textbf{[図挿入：$\rho = -0.8$ の散布図]}  
  % \textbf{[図挿入：$\rho = -1$ の散布図]}

  % \vspace{1em}

  % では，相関係数はグラフのどこに現れるかというと，\textbf{それは点のばらつきの少なさ，すなわち直線へのまとまり具合}に現れている．  
  % 相関係数の絶対値が $1$ に近づくほど，点は一直線上に集まり，反対に絶対値が小さいほど，点は直線から大きくばらついて見える．  
  % 相関係数は，「直線的な関係にどれだけ近いか」という度合いを表しているともいえる．

  % ただし，注意しておきたいのは，\textbf{相関係数は直線の「傾きの大きさ」には関係しない}という点である．  
  % たとえば，2つの直線的な関係があって，片方はゆるやかな右上がり，もう片方は急な右上がりだったとしても，  
  % 相関係数は両方とも $1$ になる．

  % 傾きは，「$X$ が1増えたときに $Y$ がどのくらい変化するか」という量を表すが，  
  % 相関係数は「ずれの方向がどれだけ一致しているか」を測るものであり，数値の変化量そのものには依存しない．

  % \vspace{1em}

  % なお，相関係数が 0 に近い場合，2つの変数の間に「明確な直線的関係がない」ことを意味する．  
  % しかしこれは，「まったく関係がない」ということと同じではない．

  % たとえば，$X$ に対して $Y = X^2$ のような曲線的な関係があったとしても，  
  % $X$ の値が負から正にわたって左右対称に分布しているような場合には，  
  % $Y$ の値は平均からのずれの符号が打ち消し合うかたちになり，  
  % 相関係数は 0 に近づくことがある．

  % ただし，これはあくまで $X$ の分布に対称性があるような場合であり，  
  % たとえば $X > 0$ の範囲に限られているときには，むしろ正の相関が現れることになる．  
  % つまり，関数の形だけで相関の強さが決まるわけではなく，\textbf{変数の分布そのもの}も関係してくるのである．

  % また，次のような場合にも注意が必要である：

  % \begin{itemize}
  %   \item $Y$ の値がすべて等しく，点が\textbf{水平な直線}上に並ぶ場合．
  %   \item $X$ の値がすべて等しく，点が\textbf{垂直な直線}上に並ぶ場合．
  % \end{itemize}

  % \noindent
  % \textbf{[図挿入：傾き0の直線に並ぶ散布図（$Y=5$ など）]}  
  % \textbf{[図挿入：傾き∞の直線に並ぶ散布図（$X=3$ など）]}

  % どちらのケースでも，散布図としては「きれいな一直線」に見える．  
  % しかし，それぞれの変数の視点で見ると，$Y$ または $X$ の値がまったく変化していないことがわかる．  
  % （たとえば，$Y$ の値が一定で傾き0の直線に乗っている場合，$X$ の値をどれだけ変えても $Y$ は同じままである．  
  % 逆に，$X$ の値が一定なら，$Y$ がどれだけ変わっても $X$ は動かない．）

  % これは，どんな観測値についても $Y$（あるいは $X$）の値が一定で，平均からのずれが常にゼロであることを意味する．  
  % すなわち，\textbf{分散 $V(Y)$ または $V(X)$ が 0 である}という状態である．

  % 相関係数の定義は，
  % \[
  % \rho_{XY} = \frac{E[(X - E[X])(Y - E[Y])]}{\sqrt{V(X)} \cdot \sqrt{V(Y)}}
  % \]
  % であり，分母に $X$ や $Y$ の分散の平方根が入っているため，このような場合には\textbf{相関係数そのものが定義されない}．

  % 直感的に言えば，どちらか一方の変数にまったくばらつきがないときには，「一緒に動いているかどうか」という問いそのものが成立しないことになる．

  % \begin{remark}
  %   本書ではピアソンの相関係数を中心に扱っているが，  
  %   相関係数には他にも，変数の\textbf{値そのもの}ではなく\textbf{順位}に注目するものもある．  
  %   たとえば，スピアマンの順位相関係数やケンドールの順位相関係数といった指標があり，  
  %   これらは直線的でない関係（たとえば右上がりの曲線など）も，ある程度反映できる特徴を持っている．

  %   ただし，何を「関係」とみなすかは分析の目的によって異なる．  
  %   ピアソンの相関係数は，変数どうしの値が「どれだけ一緒に増減しているか」という，  
  %   \textbf{方向の一致}に注目するという点で，明確な意味と使いやすさを持っている．

  %   たとえば，$Y = \sqrt{X}$ のような右上がりの曲線的関係では，ピアソンの相関係数は直線的でないため低く出ることがあるが，  
  %   順位相関係数は「順序の一致」を見るため，高い値になることがある．  
  %   この違いは，実際の散布図を見比べると明確に理解できるだろう．

  %   \textbf{[図挿入：ピアソンとスピアマンの違いを示す散布図]}
  % \end{remark}



  % \subsection{独立性と相関係数}

  % 2つの変数の間に関係があるかどうかを判断する手段として，相関係数を用いる方法をこれまで見てきた．  
  % ただし，相関係数が 0 であることと，「関係がまったくない」ことは同じではない．  
  % これは，$Y = X^2$ のような曲線的な関係では方向性が左右で打ち消し合い，  
  % 相関係数は 0 に近づいてしまう一方で，$X$ と $Y$ の間には明確な関係が存在していることからもわかる．

  % このようなときにより厳密な「関係のなさ」を捉えるために用いられるのが，\textbf{独立性}という概念である．  
  % 独立であるとは，「一方の変数がどんな値をとっても，他方のふるまいには一切影響しない」ような関係を指す．  
  % 言い換えれば，「$X$ を知っても $Y$ を予測する手がかりにはならない」という状態である．

  % この定義は，数式としては2変数の確率分布を用いて表現されるが，本書ではそうした理論には立ち入らない．  
  % ここでは，「どんなに $X$ の値が変わっても，$Y$ の傾向がまったく変わらない」ことが独立であるという直感的な理解にとどめておく．

  % 独立であるならば，変数どうしが同じ方向にずれたり，逆の方向にずれたりする傾向も一切なく，  
  % \textbf{共分散は 0 となり，相関係数も 0 になる}．  
  % \textbf{つまり，独立であれば相関係数は必ず 0 になる}という一方向の関係が成り立つ．  
  % ただし，\textbf{その逆は成り立たない}．相関係数が 0 であっても，非線形な関係や他の依存構造が存在していることはありうる．

  % したがって，相関係数が 0 という事実だけから，「2つの変数の間に関係はない」と断言することはできない．  
  % 相関係数は，あくまで「線形な方向性の強さ」に焦点を当てた指標である．  
  % より広い意味での関係や依存性を扱うには，別の枠組みが必要になる．

  % \subsubsection{独立性と相関の関係}
  % \begin{remark}
  %   もし $X$ と $Y$ が独立であれば, $E[XY] = E[X]E[Y]$ が成り立つため, 
  %   \[
  %   \mathrm{cov}(X,Y) = E[XY]-E[X]E[Y] = 0.
  %   \]
  %   ただし, 共分散が0でも非線形な依存関係が存在する可能性があるため, 必ずしも完全な独立を意味しません.
  % \end{remark}

  % \begin{prop}
  %   変数 $X$ と $Y$ が独立ならば, ピアソンの相関係数 $\rho_{XY} = 0$ となります.
  % \end{prop}

  % \begin{proof}
  %   独立性から $E[XY] = E[X]E[Y]$ となるため, 
  %   \[
  %   \rho_{XY} = \frac{E[XY]-E[X]E[Y]}{\sigma_X \sigma_Y} = 0.
  %   \]
  % \end{proof}

  % \subsection{相関から何が言えて，何が言えないか}

  % 相関係数は，2つの変数の関係を数値でとらえるための便利な指標である．  
  % しかし，この数値をどう解釈するかには注意が必要であり，誤った読み取りによって誤解や誤用が生まれることもある．  
  % この節では，相関係数を使うときに特に注意すべき点を紹介する．

  % まず重要なのは，\textbf{相関があるからといって，必ずしも片方がもう片方の原因であるとは限らない}ということだ．  
  % たとえば，夏になるとアイスクリームの売上と水難事故の件数がともに増える傾向がある．  
  % この2つの間には相関が見られるかもしれないが，それは共通の要因――この場合は気温の上昇――によって同時に影響を受けているためである．  
  % このような，第3の変数が背後にあることで生じる見せかけの相関を，\textbf{擬似相関}と呼ぶ．

  % 擬似相関には，他にもさまざまな例がある：

  % \begin{itemize}
  %   \item チョコレートの消費量とノーベル賞受賞者数には国別に見ると正の相関があるとされるが，これは教育水準や生活水準などの背景要因によるものと考えられる．
  %   \item 年賀状の枚数と収入の高さに相関が見られることがあるが，これも年齢や社会的地位といった他の要因が影響している可能性が高い．
  %   \item 教会の数と犯罪件数が正の相関を示す場合があるが，そもそも人口の多さが両方に関係している可能性がある．
  % \end{itemize}

  % また，相関があっても，\textbf{因果の向き}を取り違えることもある．  
  % たとえば，「病院の数が多い地域ほど，病気の発生件数が多い」という相関が観測されたとき，  
  % それは「病院が病気を引き起こしている」というより，「病気が多い地域だから病院が多い」と見る方が自然だろう．  
  % このように，\textbf{相関があるという事実だけから因果関係を判断することはできない}．

  % さらに，\textbf{相関が見えたり見えなかったりするのは，データの見方や切り方にも大きく依存する}．  
  % たとえば，全体で見たときには相関がほとんど見られない場合でも，年齢，地域，性別など，ある条件でデータを区切って見ると，  
  % その部分集合の中では明確な相関が現れることがある．  
  % 逆に，いくつかのグループでは強い正の相関があっても，全体をまとめてしまうと，相関が弱まったり，  
  % 場合によっては反対向きに見えてしまうことすらある．  

  % これは，\textbf{グループ間の構成や割合が全体の関係に影響を与えてしまう}ことが原因であり，  
  % このような現象は，\textbf{シンプソンのパラドックス}と呼ばれることもある．  
  % では，なぜこのような逆転が起こるのだろうか？  
  % 少し極端な例で考えてみよう．

  % ある学校で，女子生徒と男子生徒のテスト結果を分析するとする．  
  % 数学の点数と理科の点数の関係を見たところ，  
  % どちらの性別のグループ内でも「数学の点が高いほど理科の点も高い」という正の相関が見られた．  
  % しかし，男女を合わせた全体で同じ分析をすると，逆に「数学の点が高いほど理科の点は低い」という負の相関が出たとする．

  % このような現象が起きる理由は，「各グループの分布の中心（平均）や大きさ（人数）」が異なっており，  
  % グループの違いが混ざることで，本来見えていた関係が“かき消された”り“反転してしまう”ためである．

  % 実際には，グループ間で得点の平均値が大きく異なるとき，その構成比（人数の偏り）が全体の傾向を左右してしまうことがある．  
  % つまり，\textbf{「誰がどれくらいいるか」という背景情報が，相関の見かけを大きく左右する}ということだ．

  % このように，「部分を見れば右上がり（正の相関）」，「全体をまとめると右下がり（負の相関）」というような逆転現象が起きることがある．  
  % だからこそ，\textbf{相関を見るときは必ずグループ分けの構造を意識しなければならない}．

  % \textbf{[図挿入：男女グループ別には右上がり，全体では右下がりの散布図]}

  % また，ごく一部のデータ（\textbf{外れ値}）が相関係数に大きな影響を与えることもある．  
  % 多くの点がバラバラに散らばっていて相関が見えないような場合でも，  
  % 極端に離れた一つの点が相関係数を大きく押し上げる（あるいは押し下げる）ことがある．  
  % 相関係数だけを見て判断せず，\textbf{必ず散布図などで全体の形を確認すること}が重要である．

  % このように，相関係数は「関係がありそうかどうか」を数値として示す手がかりにはなるが，  
  % \textbf{それだけで原因やメカニズムを語ることはできない}．  
  % 相関係数は，あくまで「注目すべき関係があるかもしれない」という出発点を与えてくれるにすぎない．

  % だからこそ，\textbf{相関があるからといって因果があるとは限らない}ということ，  
  % そして\textbf{相関がないからといって関係がまったく存在しないとは限らない}ということを，常に念頭に置いておく必要がある．

  % なお，ここで触れたような「相関が因果を意味するとは限らない」という問題に対して，  
  % より厳密に因果関係を扱うための\textbf{因果推論（causal inference）}と呼ばれる分野がある\footnote{相関だけでは「どちらが原因でどちらが結果か」は判断できないことがある．  
  % そのようなときに，原因と結果の関係をより丁寧に分析しようとする「因果推論」という考え方がある．  
  % たとえば，「この操作をすれば結果はどう変わるか？」というような問いを扱う．  
  % 本書では詳しくは立ち入らないが，「相関だけでは不十分だ」と感じたときには，  
  % このような視点があることを思い出してほしい．}.